% \input{\pConfig/"macros-local.tex"}
\typeout{ \enumLocal  \pConfig/"macros-local.tex"}

% shortcuts: simplify typing. Type it once, check it once, use it always
% \newcommand{\apzero}{\paren{1 + \bl{\alpha_{0}}}}
\newcommand{\lst}[1]		{ \left\{ #1 \right\} }
%\newcommand{\ls}[0] 		{\href{https://en.wikipedia.org/wiki/Least_squares} {Least Squares}}
\newcommand{\ftoa}[0]				{ \href{https://en.wikipedia.org/wiki/Fundamental_theorem_of_algebra}{Fundamental Theorem of Algebra} }
\newcommand{\ftolathm}[0]			{\href{https://mathworld.wolfram.com/FundamentalTheoremofLinearAlgebra.html}{\ftola}}
%\newcommand{\quartet}[0]				{ \href{https://en.wikipedia.org/wiki/Anscombe%27s_quartet}{Anscombe's data set quartet} }
\newcommand{\lsq}[0]				{\href{https://mathworld.wolfram.com/LeastSquaresFitting.html}{Least Squares}}

% adjugate matrix (matrix of cofactors)
\newcommand{\adja}{\mat{rr}{\innerSelf{x} & -\inner{\mathbf{1}}{x} \\ -\inner{\mathbf{1}}{x} & \innerSelf{\mathbf{1}}}}
\newcommand{\astarb}{\mat{c}{\inner{\mathbf{1}}{y}\\ \inner{x}{y}}}
\newcommand{\mm}{\href{https://www.wolfram.com/mathematica/?source=home}{Mathematica}}
\newcommand{\mmalpha}{\href{https://www.wolframalpha.com/?source=home}{Wolfram Alpha}}

\newcommand{\lapacknl}{\href{https://www.netlib.org/lapack/}{Linear Algebra PACKage}}
\newcommand{\lapackwiki}{\href{https://en.wikipedia.org/wiki/LAPACK}{LAPACK}}
%\newcommand{\python}{\href{https://www.python.org/}{Python}}
%\newcommand{\matlab}{\href{https://www.mathworks.com/products/matlab.html}{MATLAB}}
%\newcommand{\octave}{\href{https://octave.org/}{GNU Octave}}

\newcommand{ \qrdec }{\mathbf{QR}}	
\newcommand{\fred}{\href{https://en.wikipedia.org/wiki/Fredholm_alternative}{\fa}}

\newcommand{\mathmnormalt}[0]		{\href{https://reference.wolfram.com/language/ref/NormalDistribution.html}{\includegraphics[ width = 1.1in, trim=0cm 0cm 0cm 0.0cm, clip ]{\pGraphics/statistics/f-normal}}}
\newcommand{\mathmnormals}[0]		{\href{https://reference.wolfram.com/language/ref/NormalDistribution.html}{\includegraphics[ width = 2.0in, trim=0cm 0cm 0cm 0.0cm, clip ]{\pGraphics/statistics/f-normal}}}
\newcommand{\mathmnormal}[0]		{\href{https://reference.wolfram.com/language/ref/NormalDistribution.html}{\includegraphics[ width = 2.5in, trim=0cm 0cm 0cm 0.25cm, clip ]{\pGraphics/statistics/f-normal}}}
\newcommand{\exponential}[0]			{\href{https://mathworld.wolfram.com/NormalDistribution.html}{e^{x} \coloneqq 1 + x + \frac{1}{2}x^{2} + \dots + \frac{1}{n!}x^{n} + \dots}}
\newcommand{\normal}[0]				{\href{https://mathworld.wolfram.com/NormalDistribution.html}{\frac{1} {\sigma\sqrt{2\pi}} \bl{\exp} \paren{-\frac{1}{2} \paren{\frac{t-\mu} {\sigma}}^{2}}}}
\newcommand{\ndf}[0]				{\href{https://en.wikipedia.org/wiki/Normal_distribution} {Normal distribution}}
\newcommand{\tfr}[0]				{\href{https://en.wikipedia.org/wiki/68-95-99.7_rule} {$68-95-99.7$ rule}}

% famous equations
\newcommand{\eqnHeat}[0]			{\href{https://mathworld.wolfram.com/HeatConductionEquation.html}{Heat Conduction}} \newcommand{\eqnLaplace}[0]			{\href{https://mathworld.wolfram.com/LaplacesEquation.html}{Laplace's Equation}}
\newcommand{\taylor}[2]{\frac{\partial #1} {\partial #2}}
\newcommand{\taylort}[2]{\frac{\partial^{2} #1} {\partial #2^{2}}}
\newcommand{\taylortt}[3]{\frac{\partial^{2} #1} {\partial #2 \, \partial #3}}

\newcommand{\relerr}[1]{\frac{\sigma_{#1}^{2}} {#1^{2}}}

\newcommand{\rdx}{\rd{x}}
\newcommand{\bly}{\bl{y}}
% called from setup-local
% \input{\pLocalSetup macros-local}

% combos
\newcommand{ \bp }[1]			{\mathbf{p}_{#1}} 
%\newcommand{\mat}[2]			{\left[\begin{array} {#1}#2 \end{array}\right]}

\newcommand{ \obspy }[0]		{\href{https://www.geophysik.uni-muenchen.de/~megies/www_obsrise/index.html\#description}{ObsPy}: \href{https://github.com/obspy}{Python Framework for Seismology}}

%\newcommand{ \yl }[1]			{{\color{yellow} {#1}}}
\newcommand{ \nedb }[0]			{\href{https://pubs.geoscienceworld.org/ssa/srl/article-abstract/81/1/12/143638/The-Nuclear-Explosion-Database-NEDB-A-New-Database}{NEDB}} 
\newcommand{ \ldeo }[0]			{\href{https://lamont.columbia.edu/}{LDEO}} 
\newcommand{ \dtic }[0]			{\href{https://discover.dtic.mil//}{DTIC}} 
\newcommand{ \osti }[0]			{\href{https://www.osti.gov/}{OSTI}}
\newcommand{ \dista }[0]			{\href{https://www.dodcui.mil/Distribution-Statements/}{Distribution A}}

\newcommand{ \gtune }[0]			{\href{https://pubs.geoscienceworld.org/ssa/srl/article-abstract/93/6/3514/617756/GTUNE-An-Assembled-Global-Seismic-Dataset-of}{Georgia Tech Underground Nuclear Explosions}}
\newcommand{ \gtuneunm }[0]		{\href{https://doi-org.libproxy.unm.edu/10.1785/0220220036}{Georgia Tech Underground Nuclear Explosions}}
\newcommand{ \mdpi }[0]			{\href{https://www.mdpi.com}{MDPI}}
\newcommand{ \sage }[0]			{\href{https://www.iris.edu/hq/sage}{SAGE}}

\newcommand{ \urlbrune }[0]		{ https://agupubs.onlinelibrary.wiley.com/doi/abs/10.1029/JB075i026p04997 } 
\newcommand{ \brune }[0]			{\href{\urlbrune}{Brune}}

\newcommand{ \urlrandall }[0]		{ https://pubs.geoscienceworld.org/ssa/bssa/article-abstract/70/1/1/101951/A-spectral-theory-for-circular-seismic-sources } 
\newcommand{ \randall }[0]		{\href{\urlrandall}{Randall}}

\newcommand{ \urlWikiSpectral }[0]	{https://en.wikipedia.org/wiki/Spectral_theory} 
\newcommand{ \spectral }[0]		{ \href{\urlWikiSpectral}{Spectral Theory} } 

\newcommand{\seqinputs}[2]		{\lst{#1_{k}, #2_{k}}_{k=1}^{m}}
\newcommand{\myline}[0]  {\arrayrulecolor{medgray}\hline}
\newcommand{\little}[0]		{\bl{\mat{c}{a_{0} \\ a_{1}}}}

%  cat  ./content/chapters/pt\ 02\ ch\ svd.tex 
% https://tex.stackexchange.com/questions/8720/overbrace-underbrace-but-with-an-arrow-instead
\NewDocumentCommand{\overarrow}{O{\U{}} O{\uparrow} m}{%
  \overset{\makebox[0pt]{\begin{tabular}{@{}c@{}}#3\\[0pt]\ensuremath{#2}\end{tabular}}}{#1}
}
\NewDocumentCommand{\underarrow}{O{\V{*}} O{\downarrow} m}{%
  \underset{\makebox[0pt]{\begin{tabular}{@{}c@{}}\ensuremath{#2}\\[0pt]#3\end{tabular}}}{#1}
}
\NewDocumentCommand{\uoverarrow}{O{\U{}_{\bymm}} O{\uparrow} m}{%
  \overset{\makebox[0pt]{\begin{tabular}{@{}c@{}}#3\\[0pt]\ensuremath{#2}\end{tabular}}}{#1}
}
\NewDocumentCommand{\voverarrow}{O{\V{}_{\bynn}} O{\uparrow} m}{%
  \overset{\makebox[0pt]{\begin{tabular}{@{}c@{}}#3\\[0pt]\ensuremath{#2}\end{tabular}}}{#1}
}
\NewDocumentCommand{\sunderarrow}{O{\sig{}_{\bymn}} O{\downarrow} m}{%
  \underset{\makebox[0pt]{\begin{tabular}{@{}c@{}}\ensuremath{#2}\\[0pt]#3\end{tabular}}}{#1}
}

\newcommand{\quartet}[0]		{\href{https://en.wikipedia.org/wiki/Anscombe\%27s_quartet}{Anscombe's data set quartet}}
\newcommand{\qquartet}[0]		{\href{https://en.wikipedia.org/wiki/Anscombe\%27s_quartet}{Anscombe's Data Set Quartet}}
%\newcommand{\anscombe}{\href{https://demonstrations.wolfram.com/AnscombeQuartet/}{\tiny$\mat{c}{a_{0} \\ a_{1}} = \mat{l}{3.0 \\ 0.50} \pm \mat{l}{1.1 \\ 0.11}$ }}
\newcommand{\anscombe}{\href{https://demonstrations.wolfram.com/AnscombeQuartet/}{\tiny$y(x) = 3.0 \pm 1.1 + (0.50 \pm 0.11 )x$ }}
\newcommand{\anscomber}{\href{https://demonstrations.wolfram.com/AnscombeQuartet/}{$y(x) = 3.0 \pm 1.1 + (0.50 \pm 0.11 )x$ }}
\newcommand{\labelx}[0]				{ \put(45.5, -12)  {$x$} }
\newcommand{\labelT}[0]				{ \put(-5, 50)  {$y$} }

\renewcommand{\soln}[0]			{\mat{c}{a_{0} \\ a_{1}}}
\renewcommand{\varx}[0]			{\frac{23}{\sqrt{2341}}}
\renewcommand{\lA}[0]			{ \mat{cc}{1 & 1  \\ 1 & 2 \\ 1 & 3 \\ 1 & 4 \\ 1 & 5 \\ 1 & 6 \\ 1 & 7 \\ 1 & 8 \\ 1 & 9} }
\renewcommand{\lB}[0]			{ \mat{cc}{1 & x_{1}  \\ 1 & x_{2} \\ 1 & x_{3} \\ 1 & x_{4} \\ 1 & x_{5} \\ 1 & x_{6} \\ 1 & x_{7} \\ 1 & x_{8} \\ 1 & x_{9}} }
%
\newcommand{\splus}[0]			{\sqrt{\half + \varx}}
\newcommand{\sminus}[0]		{\sqrt{\half - \varx}}
%
\newcommand{\corea}[0]			{ \innoo - \sigma_{-}^{2} }
\newcommand{\coreb}[0]			{ \innxx - \sigma_{+}^{2} }
%
%\newcommand{\sorea}[0]			{ \sigma_{+}^{2} - \oto }
%\newcommand{\soreb}[0]			{ \sigma_{-}^{2} - \xtx }
%
\newcommand{\sorea}[0]			{ \sigma_{+}^{2} - \innoo }
\newcommand{\soreb}[0]			{ \sigma_{-}^{2} - \innxx }
%
\newcommand{\ctone}[0]			{ \sqrt{\corea} }
\newcommand{\cttwo}[0]			{ \sqrt{\coreb} }
%
\newcommand{\stone}[0]			{ \sqrt{\sorea} }
\newcommand{\sttwo}[0]			{ \sqrt{\soreb} }
%
\renewcommand{\matrixS}[1]		{ \mat{cc}{\sigma_{+}^{#1} & 0 \\ 0 & \sigma_{-}^{#1}} }
\renewcommand{\matrixV}[0]		{ \mat{cc}{\sigma_{+} & 0 \\ 0 & \sigma_{-}} }
\renewcommand{\matrixU}[0]		{\mat{rr}{
		\sigma_{+}^{-1} \bl{\mat{r}{\cost - \phantom{2}\sint \\ \cost - 2\sint \\ \cost - 3\sint \\ \cost - 4\sint \\ \cost - 5\sint \\ \cost - 6\sint \\ \cost - 7\sint \\ \cost - 8\sint \\ \cost - 9\sint }} 
		\sigma_{-}^{-1} \bl{\mat{r}{\cost + \sint \\ 2\cost + \sint \\ 3\cost + \sint \\ 4\cost + \sint \\ 5\cost + \sint \\ 6\cost + \sint \\ 7\cost + \sint \\ 8\cost + \sint \\ 9\cost + \sint }} 
		}}
\newcommand{\matimag}{\mat{rc}{0 & i \\ -i & 0}}
\renewcommand{\lA}{\mat{rc}{1 &2 \\ -1 & 2}}
\newcommand{\ccb}{ \cellcolor{blue!20} }
\newcommand{\ccr}{ \cellcolor{red!20} }

\newcommand{\bdots}{ \bl{dots} }
\newcommand{\rdots}{ \rd{dots} }
\newcommand{\uq}{\href{https://en.wikipedia.org/wiki/Uncertainty_quantification}{Uncertainty Quantification}}
\newcommand{\uqnasa}{\href{https://ntrs.nasa.gov/api/citations/20250006412/downloads/NESC_Academy_Intro_to_UQ_2025-06-25.pdf}{Uncertainty Quantification}}
\newcommand{\swvv}{\href{https://en.wikipedia.org/wiki/Software_verification_and_validation}{Verification and Validation}}

\endinput  % -  -  -  -  -  -  -  -  -  -  -  -  -  -  -  -  -  -  -

\newcommand{\ }{\href{ }{ }}
